\chapter{Fuente de entropía: del transistor a la min-entropía}
\label{ch:ch3_entropy}

% [BORRADOR v1 -- solo teoría y simulación; cifras de hardware pendientes]
% Material de partida: docs/modelo_ero_resultados.md, docs/arquitectura.md,
% analysis/ero_model.py, analysis/jitter_estimator.py, analysis/spice_ring.py

Este capítulo construye la cadena sobre la que descansa el resto del TFM:
del jitter físico de un oscilador en anillo a una cifra defendible de
min-entropía por bit de salida, y de esa cifra a los parámetros de diseño
del muestreador. Cada paso se valida antes de usarse. Los estimadores
estadísticos se validan con datos sintéticos de jitter conocido; las cotas
analíticas se validan frente a un cálculo numérico exacto; y el banco a
nivel de transistor se valida midiendo su propia resolución antes de medir
nada con él. Todas las cifras de este capítulo son \simu{} salvo que se
indique lo contrario; la campaña de hardware se presenta en el
Capítulo~\ref{ch:ch7_results}.

\section{El oscilador en anillo elemental}
\label{sec:ero}

La fuente de entropía es un \ac{ERO}: un oscilador en anillo libre RO1
cuya salida se muestrea con un flip-flop que recibe como reloj un segundo
oscilador en anillo RO2 dividido por un entero $K_D$~\cite{baudet_2011_security,
fischer_2014_embedded}. Entre dos muestras, la fase de RO1 acumula jitter
durante $K_D$ periodos de RO2; cuando el jitter acumulado es comparable al
periodo de RO1, el bit muestreado se vuelve impredecible. Se eligió el
\ac{ERO} frente a alternativas de mayor tasa de bits (XOR de varios
anillos, muestreo coherente, anillos de efecto de transición) por una sola
razón: es la arquitectura para la que existe un modelo estocástico con una
cota de entropía demostrable y para la que se ha publicado y validado un
método embebido de medida del jitter~\cite{petura_2016_survey}. La tasa de
bits no es un requisito aquí ---el motor solo consume entropía para dar
semilla a un generador determinista---, mientras que una cota demostrable
sí lo es.

Una decisión de diseño que se deriva de la plataforma es que \emph{el
oscilador de la placa nunca se usa como referencia}. La fuente de reloj de
la Basys~3 es un oscilador MEMS con un \ac{PLL} fraccional cuya hoja de
características especifica un jitter ciclo a ciclo de hasta 95~ps
\est{}~\cite{digilent_2014_basys3rm}, entre uno y dos órdenes de magnitud
por encima del jitter por periodo de un oscilador en anillo en un proceso
de 28~nm (2--4~ps \est{}~\cite{petura_2016_survey}). Muestrear RO1 con ese
reloj mediría el ruido de fase del PLL, una componente global y
potencialmente inducible, en lugar del ruido térmico del anillo. Por eso
RO2 muestrea a RO1, y el reloj de sistema de 100~MHz solo gobierna la
lógica síncrona.

\section{Modelo estocástico y factor de calidad}
\label{sec:model}

Siguiendo a Baudet \emph{et al.}~\cite{baudet_2011_security}, la fase
$\varphi(t)$ de RO1, medida en periodos, se modela como un proceso de
Wiener con deriva $\mu$ y volatilidad $\sigma^2$. Entre dos muestras
separadas $\Delta t$, la fase avanza en promedio $\nu = \mu\,\Delta t$
periodos y acumula una varianza
\begin{equation}
  Q = \sigma^2\,\Delta t ,
  \label{eq:Q}
\end{equation}
el \emph{factor de calidad}. El bit de salida es $s = 1$ cuando la fase
fraccionaria está en $[1/2, 1)$. Condicionado a la fase anterior
$\varphi_0$, el bit sigue una distribución de Bernoulli con
\begin{equation}
  P(s=1 \mid \varphi_0) = \sum_{k\in\mathbb{Z}}
  \left[\Phi\!\left(\tfrac{k+1-\varphi_0-\nu}{\sqrt{Q}}\right) -
        \Phi\!\left(\tfrac{k+\tfrac12-\varphi_0-\nu}{\sqrt{Q}}\right)\right],
  \label{eq:puno}
\end{equation}
donde $\Phi$ es la función de distribución normal estándar; se trata de la
gaussiana envuelta evaluada de forma exacta. Su primer armónico da el sesgo
máximo $\varepsilon_{\max} \approx (2/\pi)\,e^{-2\pi^2 Q}$, a partir del
cual el Corolario~1 de Baudet proporciona la conocida expresión
\begin{equation}
  H(s \mid \varphi_0) = 1 - \frac{4}{\pi^2 \ln 2}\, e^{-4\pi^2 Q}
  + O\!\left(e^{-6\pi^2 Q}\right).
  \label{eq:baudet14}
\end{equation}

\section{Entropía condicional exacta frente a la cota asintótica}
\label{sec:exact}

La ecuación~\eqref{eq:baudet14} se cita habitualmente como una \emph{cota
inferior} de la entropía por bit. Leída con atención, la desigualdad del
corolario solo afirma que la tasa de entropía de la secuencia es al menos
la entropía de un bit condicionada a la fase anterior exacta; el miembro
derecho es un desarrollo asintótico \emph{de} esa entropía condicional, no
una cota de ella. Desarrollando la entropía binaria en función del sesgo,
$h_2(\tfrac12+\varepsilon) = 1 - \tfrac{1}{\ln 2}\sum_{n\ge1}
\tfrac{(2\varepsilon)^{2n}}{2n(2n-1)}$, y promediando sobre una $\varphi_0$
uniforme con $\varepsilon = A\sin(2\pi\varphi)$,
$A = (2/\pi)e^{-2\pi^2 Q}$, se obtiene
\begin{equation}
  H(s\mid\varphi_0) = 1 - \frac{4}{\pi^2\ln 2}\,e^{-4\pi^2 Q}
                        - \frac{8}{\pi^4\ln 2}\,e^{-8\pi^2 Q} - \cdots
  \label{eq:orden2}
\end{equation}
Todos los términos omitidos son negativos, así que \emph{cualquier}
truncamiento sobrestima la entropía; el término siguiente decae como
$e^{-8\pi^2 Q}$, no como $e^{-6\pi^2 Q}$. La Tabla~\ref{tab:exact}
compara el valor exacto, calculado a partir de \eqref{eq:puno} sobre una
malla de 8192 fases, con los desarrollos de primer y segundo orden. Añadir
el término de segundo orden reduce el error de \eqref{eq:baudet14} en un
factor de 95 para $Q=0{,}10$ y de 4970 para $Q=0{,}20$, lo que confirma la
deducción.

\begin{table}[htbp]
\centering
\caption{Entropía de Shannon condicional por bit: cálculo exacto con la
gaussiana envuelta frente a desarrollos truncados \simu{}.}
\label{tab:exact}
\begin{tabular}{@{}lcccc@{}}
\toprule
$Q$ & Exacta & Orden 2 & Orden 1 (Ec.~\ref{eq:baudet14}) & Error del orden 1 \\
\midrule
0,02 & 0,7013194 & 0,7100952 & 0,7345213 & $+3{,}3\times10^{-2}$ \\
0,05 & 0,9163001 & 0,9164920 & 0,9187783 & $+2{,}5\times10^{-3}$ \\
0,10 & 0,9886728 & 0,9886733 & 0,9887174 & $+4{,}5\times10^{-5}$ \\
0,15 & 0,9984319 & 0,9984319 & 0,9984327 & $+8{,}5\times10^{-7}$ \\
0,20 & 0,9997823 & 0,9997823 & 0,9997823 & $+1{,}6\times10^{-8}$ \\
\bottomrule
\end{tabular}
\end{table}

En la región de funcionamiento ($Q\approx0{,}2$) la discrepancia es de
$10^{-8}$ e irrelevante; para un generador infradimensionado
($Q\approx0{,}02$) la fórmula truncada atribuye un 3,3\,\% de entropía que
no existe. Por ello, este TFM usa el cálculo exacto como valor conservador
y cita \eqref{eq:baudet14} como lo que es: una aproximación asintótica
excelente para $Q \gtrsim 0{,}1$. La misma herramienta numérica muestra que
la min-entropía obtenida a partir del sesgo del primer armónico,
$H_\infty = -\log_2(\tfrac12 + \varepsilon_{\max})$, coincide con el valor
exacto del peor caso con un error inferior a $10^{-4}$~bit en
$Q\in[0{,}05;\,0{,}30]$.

\subsection{Objetivo de diseño}

La versión~3.0 de AIS~20/31~\cite{peter_2024_ais31} admite dos criterios
alternativos para la clase PTG.2: entropía de Shannon $\ge 0{,}9998$ o
min-entropía $\ge 0{,}98$ por bit interno. Resolver ambos por bisección
sobre las entropías exactas da $Q \ge 0{,}2022$ y $Q \ge 0{,}2286$,
respectivamente. El criterio de min-entropía es el más exigente y es
también la magnitud con la que trabaja NIST SP~800-90B, por lo que es el
que se adopta aquí: \textbf{el objetivo de diseño es $Q \ge 0{,}23$, con
un punto de consigna de $Q = 0{,}30$} ($H_\infty = 0{,}9951$), con un
margen que cubre la reducción del jitter térmico a baja temperatura, que
es la prueba adversaria estándar de Fischer y
Lubicz~\cite{fischer_2014_embedded}.

\section{Estimador embebido del jitter y sus dos correcciones}
\label{sec:estimator}

Para elegir $K_D$ hay que conocer $\sigma^2$, y el método de Fischer y
Lubicz~\cite{fischer_2014_embedded} la mide a partir de los propios bits:
con $K_D = 1$, la fracción $c$ de pares $(b_j, b_{j+M})$ que difieren
dentro de un bloque de $N$ pares estima la fase acumulada a lo largo de
$M$ muestras, y su varianza entre $K$ bloques crece como
$4 M Q_{\mathrm{bruto}}$. Implementado al pie de la letra, con un ajuste a
una recta, el estimador \emph{subestima} $Q_{\mathrm{bruto}}$ entre un
13\,\% y un 38\,\% según el nivel de jitter. Se identificaron dos causas
sobre datos sintéticos con $Q$ conocido, y se corrigieron:

\begin{enumerate}
\item \textbf{Efecto de vértice.} $c$ no es lineal en la fase acumulada;
es una onda \emph{triangular} de periodo uno y pendiente $\pm2$. Cuando la
fase determinista $M\nu \bmod 1$ cae cerca de un vértice, la distribución
queda a ambos lados de él y la varianza se desploma ---hasta 0,38 de su
valor nominal en un caso medido---. El vértice se detecta a partir de la
propia media medida ($\bar c \approx 0$ o $1$ es un vértice;
$\bar c \approx \tfrac12$, el centro de la rampa), de modo que las
distancias $M$ con $\bar c$ fuera de $[0{,}15;\,0{,}85]$, o para las que
$3\sqrt{MQ}$ alcanza el vértice, se descartan de forma iterativa.
\item \textbf{Ruido de estimación no constante.} Estimar una fracción a
partir de $N$ pares añade una varianza proporcional a $\bar c(1-\bar c)$,
que cambia con $M$. Ajustarla como una ordenada en el origen constante
sesga la pendiente; en su lugar se ajusta como un segundo regresor,
$\mathrm{Var}(c) = Q\,(4M) + \kappa\,\bar c(1-\bar c)/N$.
\end{enumerate}

Tras ambas correcciones, el estimador recupera el $Q_{\mathrm{bruto}}$
inyectado dentro de un $\pm2$\,\% en un rango de $150\times$
(Tabla~\ref{tab:estimator}), con un sesgo de $-0{,}6$\,\% y una dispersión
de 4,0\,\% sobre ocho semillas independientes para $K=20\,000$ bloques. El
exponente de acumulación ajustado es $1{,}009$ para un simulador sin ruido
de parpadeo; medido en hardware, el mismo exponente separa la componente
térmica ($b=1$, que cuenta como entropía) del ruido de parpadeo y el
jitter global ($b=2$, que no cuenta).

\begin{table}[htbp]
\centering
\caption{Recuperación de un $Q_{\mathrm{bruto}}$ conocido mediante el
estimador embebido corregido ($K=20\,000$ bloques de $N=100$ pares)
\simu{}.}
\label{tab:estimator}
\begin{tabular}{@{}cccc@{}}
\toprule
$Q_{\mathrm{bruto}}$ inyectado & Estimado & Error & Valores de $M$ usados \\
\midrule
$2{,}0\times10^{-7}$ & $1{,}979\times10^{-7}$ & $-1{,}07$\,\% & 19 \\
$5{,}0\times10^{-7}$ & $5{,}060\times10^{-7}$ & $+1{,}21$\,\% & 18 \\
$2{,}0\times10^{-6}$ & $2{,}037\times10^{-6}$ & $+1{,}86$\,\% & 16 \\
$8{,}0\times10^{-6}$ & $8{,}039\times10^{-6}$ & $+0{,}48$\,\% & 9 \\
$3{,}0\times10^{-5}$ & $2{,}977\times10^{-5}$ & $-0{,}78$\,\% & 5 \\
\bottomrule
\end{tabular}
\end{table}

Una consecuencia práctica para el hardware es que los bloques no tienen por
qué ser contiguos: cada uno requiere solo $N+M$ bits consecutivos, y los
bloques independientes dan una estimación algo mejor (error de
$-0{,}30$\,\% con 20\,000 fragmentos independientes). Por eso el búfer de
captura de la FPGA contiene un único fragmento de unos pocos kilobits, que
se llena a la velocidad del anillo y se lee despacio, en lugar de megabits
de muestras contiguas.

\section{Banco a nivel de transistor y su resolución}
\label{sec:spice}

La fuente de entropía es el único bloque del motor ligado a la tecnología,
y en un circuito integrado se diseña a nivel de transistor. Se construyó
un banco SPICE en torno a un anillo de cinco etapas con una puerta NAND de
habilitación, con todos los parámetros dependientes de la tecnología
(modelos de dispositivo, anchuras, longitud mínima, alimentación)
confinados en una cabecera, de modo que pasar de los modelos genéricos
provisionales usados aquí al \ac{PDK} de una fundición cambia un único
bloque de la netlist. El jitter se extrae de los instantes de cruce por el
umbral de la salida mediante la segunda diferencia
$D(M) = t_{k+2M} - 2t_{k+M} + t_k$, es decir, la varianza de Allan
aplicada a datos de fase. Su esperanza, $2M\sigma^2$, está libre del sesgo
que arrastra una diferencia simple sobre ventanas solapadas (medido en un
5\,\%), y es el método que usan Benea \emph{et al.} para caracterizar
anillos tanto en Artix-7 como en silicio de 28~nm~\cite{benea_2024_flicker}.
Las estimaciones a distintos $M$ se combinan con pesos de varianza inversa
proporcionales al número de ventanas independientes; sin esa ponderación,
la dispersión se quedaba en el 15\,\% con independencia de la longitud
simulada.

El banco mide primero su propia resolución: con todas las fuentes de ruido
desactivadas, el jitter residual es el ruido numérico del simulador. El
primer intento dio 0,33~ps, del mismo orden que el jitter que se quiere
medir; se debía por completo al transitorio de arranque del oscilador, y
descartar los primeros veinte ciclos lleva el suelo de resolución a
\textbf{1,38~fs}, tres órdenes de magnitud por debajo de la señal. Solo
entonces se inyectó ruido. El jitter es proporcional a la raíz cuadrada de
la densidad espectral inyectada ($\times2{,}01$ medido para $\times4$ en
densidad, frente a $\times2{,}00$ esperado) hasta el punto en que la
corriente de ruido eficaz alcanza aproximadamente el 80\,\% de la
corriente de conducción de la etapa; ahí deja de valer el régimen de
pequeña señal, el periodo medio se desplaza y la ley de escala falla por
un factor de cinco. De ello se deduce un criterio de validez para toda
campaña, incluidas las que se hagan sobre un \ac{PDK} real: mantener el
ruido inyectado por debajo de $\sim40$\,\% de la corriente de la etapa y
comprobar que el periodo medio no se mueve al activar el ruido. Con los
modelos provisionales el anillo oscila a 4,3~GHz, lo que no es físico; las
cifras que entran en este TFM proceden de modelos de dispositivo reales.

Los primeros modelos de ese tipo son los del \ac{PDK} abierto IHP SG13G2
(130~nm, núcleo de 1,2~V, modelo compacto PSP~103.6, licencia
Apache-2.0)~\cite{ihp_openpdk_repo}, que pueden publicarse junto con todo
lo demás. Hay dos diferencias respecto al banco con modelos provisionales.
Primero, los modelos están escritos en Verilog-A y hay que compilarlos a
la interfaz OSDI; esto obligó a compilar ngspice~47 y el compilador
OpenVAF-Reloaded en espacio de usuario, ya que ningún canal de paquetes
ofrecía versiones lo bastante recientes. Segundo, y más importante, el
ruido ya no se elige: ngspice no inyecta el ruido de los dispositivos
durante un análisis transitorio, así que la densidad de las fuentes de
corriente se toma de un análisis de ruido en pequeña señal de los mismos
transistores, un inversor polarizado en su punto de conmutación con la
entrada unida a la salida, donde ambos dispositivos conducen y el ruido
desplaza el instante de cruce. El ruido referido a la entrada de ese nodo
es exactamente la corriente que debe reproducir la fuente del transitorio.
Su espectro es un término $1/f$ puro más un suelo de ruido blanco de
$4{,}13\times10^{-24}~\mathrm{A^2/Hz}$, plano de 100~kHz a 10~GHz, con la
frecuencia de esquina del $1/f$ en 19,8~MHz y un factor de ruido efectivo
$\gamma = 1{,}04$, por encima del 2/3 de canal largo, como se espera en un
dispositivo corto. Antes de usarla, la parte $1/f$ de la fuente de ruido
de ngspice se calibró frente a una resistencia conocida: su densidad es
$\mathrm{NAMP}^2/(\pi f)$ con un margen del 0,4\,\%. Los resultados están
en la Sección~\ref{sec:ihp_ring}.

\section{Objetivo de diseño y selección del divisor}
\label{sec:kd}

Con $\sigma/T$ medido por periodo, $Q_{\mathrm{bruto}} = (\sigma/T)^2$ y
el divisor se obtiene de $K_D = \lceil Q_{\mathrm{objetivo}} /
Q_{\mathrm{bruto}} \rceil$. Para el rango de jitter esperado en un proceso
de 28~nm, el divisor queda entre $3\times10^4$ y $5\times10^5$, y la tasa
de bits por anillo entre 1 y 20~kbit/s. El modelo reproduce el punto de
diseño publicado por Fischer y Lubicz ($K_D \approx 430\,000$ para
$\sigma = 5{,}01$~ps sobre $T = 7{,}81$~ns)~\cite{fischer_2014_embedded}.
Estas tasas bastan aquí: dar semilla al generador determinista con
256~bits de entropía completa requiere unos 512 bits brutos con
$H_\infty \approx 0{,}98$, es decir, entre 0,03 y 0,5~s por semilla, y el
motor solo emite una semilla en el arranque y en los intervalos de
resembrado. Si alguna vez hiciera falta más tasa de bits, el camino limpio
son varios \acp{ERO} independientes, cada uno con su propia cota, en lugar
de un XOR de varios anillos cuya independencia no puede
demostrarse~\cite{bochard_2010_true}.
